\section{The actual discriminant cover in the even-quadric cycle gate}

\subsection{BR-15-ORBIT-RETRACT: recovering finite Chow retractions from an orbit isomorphism}
Let $\mathcal C$ be an additive category with an invertible additive
autoequivalence $M\mapsto M(1)$. Its orbit category has
\[
 \operatorname{Hom}_{\mathcal C/(1)}(M,N)
 =\bigoplus_{j\in\mathbf Z}\operatorname{Hom}_{\mathcal C}(M,N(j)),
\]
with the usual shifted composition. If $M$ and $N$ are isomorphic in
this orbit category, then $M$ is a retract in $\mathcal C$ of a finite
direct sum of twists of $N$, and conversely $N$ is a retract of a finite
direct sum of twists of $M$. Consequently $M$ and $N$ belong to the same
full replete additive subcategories closed under retracts and all twists.

\textbf{Proof.} Write an orbit isomorphism and its inverse as finite
families $u=(u_j)$ and $v=(v_k)$, with
$u_j:M\to N(j)$ and $v_k:N\to M(k)$. The component of twist zero in
$vu=1_M$ is the identity of actual morphisms in $\mathcal C$
\[
 \sum_j v_{-j}(j)u_j=1_M. \tag{BR-OD.1}
\]
Choose the finite set $J$ of indices with $u_j\ne0$, and put
$P=\bigoplus_{j\in J}N(j)$. The maps
\[
 a=(u_j)_{j\in J}:M\longrightarrow P,\qquad
 b=(v_{-j}(j))_{j\in J}:P\longrightarrow M
\]
satisfy $ba=1_M$ by (BR-OD.1). Thus $M$ is a retract of $P$. Interchange
$u$ and $v$ for the other retraction. Closure under finite sums, twists
and retracts proves the final assertion. All sums here are finite
because the orbit Hom is a direct sum. No conservativity of a realization,
nilpotence conjecture, or lifting of a numerical idempotent is used.
$\square$

\subsection{BR-15-QUADRIC-DISCRIMINANT: replacing the unspecified surface projector}
Let $k$ be a field of characteristic different from two. Let $B$ be a
smooth projective $k$-scheme, and let $f:Q\to B$ be the flat, generically
smooth quadric
fibration associated with a line-bundle-valued quadratic form on a
vector bundle of even rank $r+2$, where $r\ge2$ is even. Assume $Q$ is
smooth, every geometric fibre has corank at most one, and the
discriminant divisor $D\subset B$ is smooth. We allow an empty
discriminant divisor. Define the discriminant double cover
\[
 T=\operatorname{Spec}_B Z(\operatorname{Cl}_0(f)). \tag{BR-OD.2}
\]
This is a smooth projective $k$-scheme; it need not be geometrically
connected. For every full replete additive subcategory $\mathcal A$ of
rational Chow motives over $k$ that is closed under direct summands and
all Tate twists,
\[
 h(Q)\in\mathcal A
 \quad\Longleftrightarrow\quad
 h(B)\in\mathcal A\ \hbox{and}\ h(T)\in\mathcal A. \tag{BR-OD.3}
\]
In particular this holds for the Chow category generated by abelian
varieties and Artin motives. After algebraic closure, it also holds for
the subcategory of Tate motives and for the subcategory generated by
supersingular abelian varieties. This identifies a concrete geometric
test in the even-relative-dimensional case of CC-XC18.
The same equivalence holds for membership in a full additive,
summand- and twist-closed category of rational $\ell$-adic homological
motives, for any invertible prime $\ell$. In particular it identifies
the original homological abelian--Artin forward-cycle gate with the
two gates for $B$ and the specified cover $T$.

\textbf{Exact published inputs.} We use rational coefficients throughout.
Bernardara--Marcolli--Tabuada (2018, Proposition 4.1, pp.~306--307,
diagram (4.2)) give a fully faithful additive functor
\[
 R:\operatorname{Chow}(k)_{\mathbf Q}/(-\otimes\mathbf Q(1))
   \longrightarrow\operatorname{NChow}(k)_{\mathbf Q}.
 \tag{BR-OD.4}
\]
It takes the orbit of $h(V)$ to the noncommutative Chow motive
$U(\operatorname{perf}_{\rm dg}V)$ of a smooth projective scheme $V$;
see their p.~307, proof of Theorem 1.1. Their equation (6.1) and the
proof of Theorem 1.2(ii), p.~308, give in this projective setting
\[
 U(Q)\simeq U(T)\oplus U(B)^{\oplus r}.
 \tag{BR-OD.5}
\]
Here the first decomposition comes from the Clifford component of the
quadric fibration; the even Clifford algebra becomes Azumaya on $T$;
and its rational noncommutative motive agrees with that of $T$.
The Azumaya step is Tabuada--Van den Bergh (2013, Theorem 2.1,
arXiv:1307.7946v2, p.~3): the required ranks are invertible in
$\mathbf Q$.

The field hypothesis deserves an explicit source qualification.
Kuznetsov's original 2008 paper assumes an algebraically closed field
of characteristic zero in its \S2, p.~1343. The general-field
quadric semiorthogonal decomposition used here is stated and justified
by Auel--Bernardara--Bolognesi (2014, Theorem 2.2.1 and its proof,
p.~270), which expressly removes that restriction. Tabuada (2020,
Proposition 15(i), p.~2346, complete proof on that page) gives the same
decomposition after inverting two and describes the smooth discriminant
cover, under the explicit simple-degeneration hypotheses above. We do
not transfer a characteristic-zero-only theorem to finite fields by
analogy.

\textbf{Complete deduction.} Put $N=h(T)\oplus h(B)^{\oplus r}$.
Equation (BR-OD.5) and full faithfulness of (BR-OD.4) lift to an
isomorphism $h(Q)\simeq N$ in the Tate orbit category: lift each
direction by fullness, and use faithfulness to lift the two inverse
identities. BR-15-ORBIT-RETRACT now gives actual Chow-motive
retractions
\[
 h(Q)\ \mid\ \bigoplus_{j\in J}N(j),\qquad
 N\ \mid\ \bigoplus_{k\in K}h(Q)(k) \tag{BR-OD.6}
\]
for finite sets $J,K$, where $M\mid P$ denotes a specified split
inclusion. If $h(B)$ and $h(T)$ belong to $\mathcal A$, the first
retraction puts $h(Q)$ in $\mathcal A$. Conversely the second puts
$N$ in $\mathcal A$, and projection onto its displayed summands gives
both $h(B),h(T)\in\mathcal A$; here $r\ge2$ ensures that a copy of
$h(B)$ is present. This proves (BR-OD.3). These are actual rational
Chow correspondences with identities modulo rational equivalence.
The original degree-zero-component argument does not specify a
particular graded direct-sum formula for $h(Q)$, and none is needed.

All the hypotheses, the Clifford centre, and the quoted decompositions
commute with field extension. Apply the same proof over $\bar k$ to
obtain the last two assertions. The subcategories in question are
closed under finite sums, summands, and twists by their definitions.
For the homological assertion, let $\mathcal A$ be the full preimage
of the specified homological subcategory under the rational Chow
realization functor. That preimage has all the closure properties used
in (BR-OD.3), since a functor preserves a displayed retraction. Thus
(BR-OD.3) applies to it. This argument assumes no equality between
numerical and homological equivalence for $B$, $T$, or $Q$.
$\square$

\subsection{BR-15-QUADRIC-TATE: finite-field cycle consequences and a concrete K3 case}
Suppose now that $k=\mathbf F_q$ has odd characteristic and the
preceding hypotheses hold. If $B_{\bar k}$ and $T_{\bar k}$ have Tate
Chow motives, then $h(Q)$ is Artin--Tate over $k$. On all finite
products of such total spaces, and after every finite base extension,
rational, homological and numerical equivalence coincide; the
rational cycle-class map tensored with $\mathbf Q_\ell$ is an
isomorphism onto the Frobenius-invariant even cohomology; and odd
cohomology vanishes. If the two geometric generators are instead of
supersingular abelian Chow type, the corresponding conclusions of
CC-XC22 hold, including semisimple Frobenius, without an assertion that
odd cohomology vanishes. In either case the fixed-target forward-cycle
search terminates.

\textbf{Proof.} BR-15-QUADRIC-DISCRIMINANT puts the geometric motive of
$Q$ in the same stated subcategory as its two generators. Apply
CC-XC20 or CC-XC22, respectively, whose proofs give actual finite-field
Chow splittings and all-product comparisons. This proves every
assertion, without the all-abelian 2026 Tate input in the supersingular
cases. $\square$

For example, take $B=\mathbf P^2_{\mathbf F_q}$ and suppose the
actual smooth discriminant double cover $T$ is a geometrically
supersingular K3 surface. The projective bundle decomposition gives
$h(B)=\mathbf1\oplus\mathbf1(-1)\oplus\mathbf1(-2)$, and the
published supersingular K3 result used in CC-XC21 makes $h(T_{\bar k})$
Tate. Therefore all the preceding Tate conclusions hold for $Q$.
This condition refers to the explicitly defined cover (BR-OD.2),
rather than to an unidentified projector of some auxiliary surface.

There is also a precise birational corollary. Let $Y/\mathbf F_q$ be a
smooth projective variety for which a smooth projective $Q$ satisfying
this Tate criterion admits a proper birational morphism $\beta:Q\to Y$.
Then the same conclusions hold for $Y$ and its finite products.
Indeed, the graph and transpose graph give actual Chow correspondences
$\beta^*,\beta_*$, and $\beta_*\beta^*=1_{h(Y)}$: the graph
composition is supported on the diagonal, and its generic multiplicity
is $\deg\beta=1$. Equivalently, on every product with a smooth
projective test scheme this is the projection formula and the
pushforward of the fundamental cycle. Thus $h(Y)$ is a Chow summand of
$h(Q)$, also after base change, and CC-XC20 applies.

In particular the conclusion applies to a smooth cubic fourfold
containing a plane whenever its blowup along that plane has the
displayed flat quadric-surface fibration over $\mathbf P^2$, smooth
simple discriminant, and geometrically supersingular K3 discriminant
cover. The blowup is smooth and its morphism to the fourfold is
birational. No assertion that every cubic with a plane has this
discriminant geometry is made.

\subsection{Exact references and scope}
\begin{itemize}
\item Auel, A., Bernardara, M., \& Bolognesi, M. (2014).
Fibrations in complete intersections of quadrics, Clifford algebras,
derived categories, and rationality problems.
\emph{Journal de Math\'ematiques Pures et Appliqu\'ees, 102}(1),
249--291. \url{https://doi.org/10.1016/j.matpur.2013.11.009}.
Theorem 2.2.1 and general-field proof, p.~270; Proposition 1.2.5,
p.~257, for smooth discriminant and simple degeneration.
\item Bernardara, M., Marcolli, M., \& Tabuada, G. (2018).
Some remarks concerning Voevodsky's nilpotence conjecture.
\emph{Journal f\"ur die reine und angewandte Mathematik, 738},
299--312. \url{https://doi.org/10.1515/crelle-2015-0068}.
Orbit category definition, p.~305; Proposition 4.1, pp.~306--307;
equation (6.1) and proof of Theorem 1.2(ii), p.~308.
\item Tabuada, G. (2020). Schur-finiteness (and Bass-finiteness)
conjecture for quadric fibrations and families of sextic du Val del
Pezzo surfaces. \emph{Documenta Mathematica, 25}, 2339--2354.
\url{https://doi.org/10.4171/DM/800}. Theorem 1, p.~2340;
Proposition 15(i), complete proof and discriminant-cover definition,
p.~2346.
\item Tabuada, G., \& Van den Bergh, M. (2013).
\emph{Noncommutative motives of Azumaya algebras}
(arXiv:1307.7946v2). \url{https://arxiv.org/pdf/1307.7946v2}.
Theorem 2.1, p.~3; proof \S8, pp.~14--18. The numbered statement is
used as a cited input; no new proof of Azumaya invariance is claimed.
\end{itemize}

The known decompositions and fully faithful functor are cited inputs.
The finite orbit-retraction argument and its consequences are proved
above. The corank, smoothness, line-bundle quadratic presentation and
odd-characteristic conditions have not been dropped. Arbitrary even
quadric fibrations remain under the broader, less explicit CC-XC18
reduction when these additional hypotheses are unavailable.
